\scanpage{39}
D'où on conclut, utilisant l'équivalence de la catégorie des rev. étale de $\unsure{Y|X}$ et des rev. étale de $X_{/Y}$, le

\uline{Corollaire} Si $(LG)$ vaut, alors le foncteur $X'\leadsto X'\times_X Y$ des rev. étales de $X$ dans les rev. étales de $Y$ est pleinement fidèle, \deleted{Si $X$ $(LG)$ strict vaut, pour} et kif kif en remplaçant $Y$ par un voisinage ouvert $U$ quelc de $Y$. Si $(LG)$ strict vaut, \deleted{alors} pour tout rev étale $Y'$ de $Y$, il existe un \deleted{rev étale $\unsure{X'|U'}$ d'un} voisinage ouvert $U$ de $Y$, \deleted{et} un rev. étale $U'$ de $U$, \deleted{tel que} \inserted{\unsure{en|et} \unsure{plus|bien}} un $Y$-isom $U'\times_U Y\simeq Y'$. (Si $U$ est choisi, $U'$ est donc déterminé à isom unique près).

\uline{Remarque}. Nous verrons plus loin des conditions, liées au ``Th de pureté'' permettant de conclure que $\pi_1(Y)\longrightarrow\pi_1(X)$ est \deleted{surjectif}\inserted{bijectif}. Kif-Kif pour Picard ----
% archive sheet number: 38
